\section{Related Work}
\label{sec:related}


Here, we discuss related work on weight decay and model plasticity and how our work contributes new insights.

\textbf{Weight decay in language model training.}
Weight decay is a standard hyperparameter in language model training and is commonly implemented in conjunction with adaptive optimizers such as AdamW~\citep{loshchilov2019decoupled, brown2020language, grattafiori2024llama, olmo2024olmo2, liu2024deepseekv3}.
Beyond its classical role in regularization and generalization~\citep{krogh1991simple,zhang2018three,loshchilov2019decoupled,zhou2024towards}, weight decay has also been shown to play other roles in language model training, such as improving optimization and training stability~\citep{andriushchenko2023whywd}, shaping the learning rate~\citep{li2020reconciling,kosson2024rotational,kosson2025weight}, controlling the effective step size~\citep{wen2025hyperball}, inducing low-rank attention layers~\citep{kobayashi2024wdlowrank}, and increasing forgetting of contaminated benchmark questions~\citep{bordt2025datacontam}. \citet{wang2024set} show that the weights of AdamW can be understood as an exponential moving average, and that the weight decay hyperparameter plays a critical role in controlling its time scale. \citet{bergsma2505powerlines} study how to set weight decay to minimize the pretraining loss of language models, finding that lower weight decay improves pretraining loss in the overtrained (high TPP ratio) regime. \citet{kim2025ptunderinfcompute} show that larger weight decay improves pretraining loss in the multi-epoch setting. In contrast to previous work which primarily focuses on weight decay's effects on the pretrained model, this paper examines how weight decay during pretraining shapes model plasticity. %To our knowledge, this paper is the first work to investigate the role of weight decay in language model plasticity.




\textbf{Plasticity of deep learning models.}
Model plasticity has previously been studied in the contexts of continual learning, transfer learning, and reinforcement learning, settings in which models often undergo multiple rounds of training~\citep{dohare2024loss, klein2024plasticity,coetzer2025restoring}.
Prior works have demonstrated that image models lose plasticity when subjected to additional rounds of training on new data, leading to a decreased ability to learn this new data~\citep{dohare2024loss,lyle2023understanding,klein2024plasticity}.
Various approaches have been developed to improve model plasticity, including shrinking and perturbing model weights at the start of each training round~\citep{ash2020shrinkandperturb}, identifying and re-initializing less-useful model weights during training~\citep{dohare2024loss}, pushing weights towards initialization weights during training~\citep{kumar2023maintaining}, and learning per-connection plasticity strengths among neuron pairs~\citep{miconi2018differentiable}.
While previous studies have examined how active forgetting and tokenization~\citep{chen2023improving,abagyan2025onetokenizer} affect language model plasticity, research on language model plasticity remains underdeveloped. In contrast to these works, this paper investigates the role of weight decay, a standard hyperparameter for language model training, on language model plasticity.


